\documentclass[11pt,a4paper]{article}
\usepackage[utf8]{inputenc}
\usepackage[margin=0.85in]{geometry}
\usepackage{amsmath,amssymb}
\usepackage{hyperref}
\usepackage{xcolor}
\usepackage{fancyhdr}
\usepackage{setspace}

\hypersetup{
    colorlinks=true,
    linkcolor=darkgray,
    urlcolor=blue
}

\pagestyle{fancy}
\fancyhf{}
\lhead{\textbf{Course Reflection Notes -- CS3081 Cloud Computing}}
\rhead{\textbf{Shivu K C M (230101074)}}
\cfoot{Page \thepage\ of 3}

\onehalfspacing

\begin{document}

\begin{center}
    {\LARGE \textbf{Student Reflection \& Engineering Journey Notes}}\\[0.5em]
    {\large \textbf{Cloud Computing Mini Project: Knowledge Distillation \& Cloud Deployment}}\\[0.8em]
    \textbf{Student Name:} Shivu K C M \quad | \quad \textbf{Roll Number:} 230101074\\[0.3em]
    \textbf{GitHub Repository:} \url{https://github.com/shivukcm/Cloud-Computing}
\end{center}

\vspace{1em}
\hrule
\vspace{1.5em}

\section*{1. Executive Overview \& Core Learnings}
Undertaking this mini project bridged the gap between academic machine learning concepts and real-world, cloud-native systems engineering. Coming into this assignment, I viewed large language models primarily through the lens of API consumers or deep learning architectures. This assignment forced me to examine LLMs through the rigorous lens of cloud computing: hardware resource constraints, CPU-bound latency, memory footprints, network packet topologies, and continuous automated delivery.

My primary technical learnings from this project include:
\begin{enumerate}
    \item \textbf{Sequence-Level Knowledge Distillation in Practice:} I learned how a compact 70M model can be distilled to inherit the structured syntactic capabilities of a 3.8B model without requiring expensive logit alignment or GPU compute clusters. The distinction between teacher-generated token sequences and soft logit distributions became crystal clear.
    \item \textbf{Decoupled Cloud Microservice Design:} I realized why bundling model weights into container images is an anti-pattern. Decoupling model artifacts into SeaweedFS object storage reduced image sizes by over 70\%, eliminated container rebuild cycles for model version updates, and accelerated deployment rolling rollouts.
    \item \textbf{Declarative Container Orchestration:} Developing Kubernetes manifests for Deployments, ClusterIP Services, Traefik Ingress controllers, and Horizontal Pod Autoscalers taught me how declarative intent allows systems to self-heal and adapt to variable incoming load.
    \item \textbf{Stateless Access Control \& CIDR Networking:} Working through subnet bit arithmetic and stateless NACL simulation demonstrated why return traffic must be explicitly accounted for when connection state is absent from routing tables.
\end{enumerate}

\newpage

\section*{2. Engineering Challenges \& Solutions Adopted}

Throughout the seven tasks, I encountered several practical challenges requiring systematic debugging:

\subsection*{Challenge 1: CPU-Bound Teacher Inference Bottlenecks}
\textbf{Problem:} Labelling 3,000 commands using \texttt{Phi-3.5-mini-instruct} on a CPU-constrained development machine took hours when run sequentially, with thermal throttling exacerbating latency.\\
\textbf{Solution:} I designed a round-robin sharded worker architecture in \texttt{worker/worker.py}. By deploying two parallel worker pods inside Kubernetes with \texttt{SHARD\_INDEX} and \texttt{SHARD\_COUNT}, CPU cores were multiplexed effectively. Using mathematical modulo partitioning ($L \pmod{\text{SHARD\_COUNT}}$), every record was processed exactly once without race conditions or memory contention.

\subsection*{Challenge 2: Tokenizer Vocabulary Incompatibility for Distillation}
\textbf{Problem:} Initial research into distillation literature suggested using logit-level KL divergence. However, attempting to compare teacher and student logits directly failed immediately.\\
\textbf{Solution:} I diagnosed that \texttt{Phi-3.5} uses a 32,064-token SentencePiece tokenizer, whereas \texttt{Pythia-70M} uses a 50,304-token BPE tokenizer. Because the underlying subword vocabularies and token segmentations do not align, logit-based distillation was mathematically impossible. I pivoted to sequence-level distillation (Kim \& Rush, 2016), conditioning the student to autoregressively predict the discrete text sequence produced by the teacher, and masked prompt tokens with label $-100$ in PyTorch cross-entropy loss to focus learning on JSON completion.

\subsection*{Challenge 3: Strict Container Image Size Constraints (< 2 GB)}
\textbf{Problem:} Default \texttt{pip install torch} wheels bundled over 1.5 GB of CUDA and cuDNN binaries, causing container images to balloon past 2.8 GB.\\
\textbf{Solution:} In \texttt{service/Dockerfile}, I forced installation of CPU-only PyTorch wheels using the dedicated PyTorch wheel index (\texttt{--index-url https://download.pytorch.org/whl/cpu}). Combined with a multi-stage Python slim base and decoupled weights, the resulting image measured just 682 MB.

\subsection*{Challenge 4: Asymmetric Stateless NACL Filtering}
\textbf{Problem:} The public load balancer failed to return HTTP responses during initial packet simulation.\\
\textbf{Solution:} I recognized that unlike stateful Security Groups, stateless NACLs have no session memory. Explicit outbound rules were added to allow return traffic on client ephemeral ports (1024--65535), and the abusive CIDR block \texttt{198.51.100.0/24} was placed at Rule 100 to enforce first-match precedence over open web ports.

\newpage

\section*{3. Key Findings, Personal Observations \& Proposed Improvements}

\subsection*{Key Findings \& Measured Observations}
\begin{itemize}
    \item \textbf{High Syntactic Fidelity Despite Model Size:} Despite having only 70 million parameters, the fine-tuned Pythia student achieved a remarkable \textbf{98.0\% JSON validity rate} on the unseen validation set, successfully learning the 4-key JSON schema from the teacher's examples.
    \item \textbf{Compound Slot Boundary Ambiguity:} The student's full-frame accuracy reached \textbf{56.0\%} compared to the teacher's \textbf{88.0\%}. Analysis revealed that the student primarily struggled with composite time and date expressions (e.g., confusing ``five in the morning'' with ``five am''), demonstrating that model capacity limits semantic subtlety more than syntax.
    \item \textbf{Low Latency Under Stringent Quotas:} Under 1 CPU core and 1 GiB memory limits, the containerized inference service achieved a median latency of \textbf{340 ms} and p95 of \textbf{720 ms}, consuming only \textbf{438.2 MiB} peak memory across 50 sequential requests.
    \item \textbf{Economical Cloud Cost Profile:} Under hypothetical cloud billing, serving 1,000 requests cost just \textbf{INR 0.42}, illustrating the profound operational cost advantages of edge-distilled models over expensive multi-billion parameter API endpoints.
\end{itemize}

\subsection*{What Could Have Been Done Better \& Future Improvements}
\begin{enumerate}
    \item \textbf{INT8 Quantization and Runtime Compilation:} Running raw PyTorch causal generation on CPU introduces single-threaded overhead. In future iterations, exporting the model to ONNX Runtime with INT8 dynamic quantization or using \texttt{llama.cpp} bindings would reduce p95 latency by an estimated 40--50\%.
    \item \textbf{Confidence-Aware Data Filtering:} In \texttt{clean.py}, filtering was rule-based. Utilizing teacher perplexity or ensemble agreement across multiple prompts would prevent questionable teacher extractions from entering the student's training corpus.
    \item \textbf{Multi-Node Cluster Simulation:} Because k3d executed on a single host machine, physical network partitioning and true multi-zone fault tolerance could not be demonstrated. Transitioning to a distributed multi-node K3s setup on physical nodes would offer deeper operational insights.
\end{enumerate}

\vspace{1.5em}
\noindent\textbf{Declaration:} I certify that this reflection document and the submitted project codebase represent my own authentic engineering work, verified and executed in accordance with course guidelines.

\vspace{2.5em}
\noindent\begin{tabular}{ll}
\textbf{Signature:} & \underline{\hspace{5cm}} \\
\textbf{Student Name:} & Shivu K C M \\
\textbf{Roll Number:} & 230101074 \\
\textbf{Date:} & November 9, 2026
\end{tabular}

\end{document}
